\documentclass{RPreport}

% Es kann passieren, dass die Befehle vom Textsatzprogramm (TexStudio) nicht erkannt werden
% Dann die cwl-Datei unter C:\Users\Benutzername\AppData\Roaming\texstudio\completion\user ablegen und TeXstudio neu starten
% Das ist optional! Normalerweise sollte die cwl-Datei automatisch generiert werden
% Das Dokument lässt sich auch kompilieren wenn die Befehle vom Textsatzprogramm nicht erkannt werden
% Unbenutzte Befehle können gelöscht werden

% Text eintragen
% --- Kopfzeile
\KopfzeileOben{z.B. Laborbericht ...}
% --- Deckblatt
\Thema{z.B. Messung von ...}
\Art{z.B. Bericht zum Laborpraktikum}
\Datum{\today} % kann man auch manuell reinschreiben
\Gruppe{1.2.3}

% Wenn Arbeit (s.u. \Arbeit):
\Matrikelnummer{1234567}
\Strasse{Musterstraße 1}
\Ort{12345 Stadt}
\EMail{name@th.de}
\Firmenname{Beispiel GmbH}

% m oder w angeben, mehrere Namen mit \newline trennen, Unterstreichen mit \underline{Name}
\Bearbeiter{m}{Roman Pfleger}
\Betreuer{w}{Betreuerin}

% Ja (j) / Nein (n)
\Arbeit{n} % Projekt- oder Abschlussarbeit: Ändert Layout und Deckblatt
\ErsteSeiteLeer{n}
\Deckblatt{j}
\ZwischenSeiteLeer{n}
\Inhaltsverzeichnis{j} % zum Aktualisieren zwei mal kompilieren
\Abbildungsverzeichnis{n}
\Tabellenverzeichnis{n}
\Autorenverzeichnis{n}
\Kopfzeile{j}
\LogoInKopfzeile{j}

% für Arbeit:
\Druckversion{n} % formatiert das Dokument im Duplexmodus

% Dateinamen ohne .tex eintragen oder leer lassen
\Titelbilddatei{} % Default ist das VT-Ohm-Logo
\Abstractdatei{Abstract}
\Symbolverzeichnisdatei{Symbole}

% .bib Datei eintragen
\addbibresource{Literatur.bib}

% Wenn Arbeit:
\Danksagungdatei{}
\Sperrvermerkdatei{}
\Firmenlogodatei{} % Wird auf dem Deckblattneben dem Ohm-Logo angezeigt

% Wenn Autorenverzeichnis:
\Autor{Max Mustermann}{Universität Musterstadt}{max@uni.de}
\Autor{Erika Beispiel}{Institut für Demo-Forschung}{erika@demo.org}

% LaTeX-intern
\title{Ein Beispiel mit Autorenverzeichnis in TeXstudio}
\author{Projektteam Demo}
\date{\today}

\begin{document}
	\InitLayout % Layout initialisieren
	\RPintro % stellt den ganzen Anfang zusammen
	% Abstract und andere Verzeichnisse werden oben eingebunden
	
	\chapter{Einleitung}

Die Messung von Durchflüssen ist ein fundamentaler Aspekt in zahlreichen Bereichen der Technik. Die präzise Erfassung von Stoffströmen spielt eine entscheidende Rolle für Effizienz, Sicherheit und Qualitätskontrolle, denn die Regelung und Bilanzierung der Stoffströme kann nur mit genauen Werten durchgeführt werden. In diesem Praktikumstermin wurden zwei Versuche, mit dem Ziel, die Durchflüsse verschiedener Medien kennenzulernen, durchgeführt. Der erste mit dem inkompressiblen Medium Wasser, der zweite mit dem kompressiblen Medium Luft. Die Vorgaben zur Versuchsdurchführung und Hintergründe dafür sind durch die Praktikumsanleitung gegeben. \cite{waermeatlas} \cite{CanJour}
	\chapter{Testtext}
\section{Theoretische Grundlagen}
\subsection{Durchflussmessung}

Die Durchflussmessung in prozesstechnischen Anlagen ist ein komplexes Thema, das von einer Vielzahl an Parametern und Rahmenbedingungen abhängt. In Rohrleitungen wird der Durchfluss maßgeblich durchFaktoren wie die spezifischen Stoffeigenschaften des Fluids, die Charakteristiken der Fördermaschine, vorhandene Höhen- und Druckunterschiede sowie die Eigenschaften der Rohrleitung selbst, wie Länge oder Rauigkeit, beeinflusst.\\
Integrierte Durchflussstelleinrichtungen und Messsysteme können den Durchfluss mehr oder weniger stark modifizieren. Die Vielfalt der möglichen Prozessmedien erfordert entsprechend differenzierte Messmethoden. Diese müssen die chemisch-physikalischen Eigenschaften wie Zusammensetzung, Leitfähigkeit und Viskosität ebenso berücksichtigen wie mögliche Aggregatzustände und Gemischzusammensetzungen. Falls die Zuverlässigkeit im Vordergrund steht, kann ein Durchflusswächter, der ein binäres Signal liefert, ausreichend sein, doch meistens ist die Genauigkeit ein wichtiges Auswahlkriterium.

\subsection{Messverfahren}
\subsubsection{Wirkdruckverfahren}

Das Wirkdruckverfahren basiert auf der Bernoulli-Gleichung~(\ref{eqn: Bernoulli}) und der Kontinuitätsgleichung~(\ref{eqn: Konti}).
\begin{align}
	\label{eqn: Bernoulli}
	p_1 + \frac{1}{2} \rho v_1^2 + \rho g h_1 &= p_2 + \frac{1}{2} \rho v_2^2 + \rho g h_2\\
	\label{eqn: Konti}
	A_1 v_1 &= A_2 v_2
\end{align}
\begin{enumerate}[]
	\item[] $p$: Druck
	\item[] $\rho$: Dichte
	\item[] $v$: Geschwindigkeit
	\item[] $g$: Erdbeschleunigung
	\item[] $h$: Höhe
	\item[] $A$: Querschnittsfläche
\end{enumerate}

Bei der Wirkdruckmessung wird ein Strömungswiderstand, i.d.R. Düsen oder Blenden nach DIN 1952, in die Strömung eingebaut, der einen Druckunterschied (Wirkdruck) zwischen beiden Seiten verursacht. Der nichtlineare Zusammenhang zwischen Volumenstrom $Q_v$ und Wirkdruck $\Delta$p wird bei Vernachlässigung der Reibung durch Gleichung~() beschrieben.
	% hier weitere Kapitel
	\include{Zusammenfassung}
	
	\printbibliography[heading=bibintoc]
	\appendix
	\chapter{Anhang}
Große Tabellen mit Messergebnissen oder Diagramme oder andere Unterlagen

	\RPoutro % stellt den ganzen Schluss zusammen
	% Autoren werden unten eingebunden

\end{document}
